% !TEX root = ../main.tex
%% Chapitre 16 — Probabilités conditionnelles
%% Sources : frontend/cours/resume/terminal/resume_probabilites_conditionnelles_terminal.txt
%%           frontend/exercices/terminal_spé_Maths/exercices_probabilite_conditionnel.txt
%% Exercice 10, question 4 : remplacée (elle reprenait la question 3 de l'exercice 8).
%% Exercices 12, 13 et 14 : problèmes inédits rédigés pour le livre.

\begin{ouverture}[conditionnelles]{16}{Probabilités conditionnelles}{Partie 3 · Probabilités}
  {Conditionnement\enspace·\enspace Arbres pondérés\enspace·\enspace Probabilités totales\enspace·\enspace
   Indépendance\enspace·\enspace Épreuves successives}
  \objectif{Calculer une probabilité conditionnelle et utiliser la relation $P(A\cap B)=P(A)\times P_A(B)$.}
  \objectif{Construire et exploiter un arbre pondéré ou un tableau à double entrée.}
  \objectif{Appliquer la formule des probabilités totales et « inverser » un conditionnement.}
  \objectif{Reconnaître et utiliser l'indépendance de deux événements.}
  \objectif{Modéliser une succession d'épreuves, y compris par une suite de probabilités.}
  \prerequis{Probabilités : événements, réunion, intersection, contraire ; indépendance de deux événements (Seconde et Première).}
  \prerequis{Dénombrement (chapitre 15) ; suites arithmético-géométriques (chapitre 1).}
\end{ouverture}

%% =====================================================================
%%  L'ESSENTIEL
%% =====================================================================
\essentiel

\essSection{1}{Probabilité conditionnelle}

\begin{definition}[Probabilité de $B$ sachant $A$]
Soit $A$ un événement de probabilité non nulle. La probabilité de $B$ \textbf{sachant} $A$ est
\[ P_A(B)=\frac{P(A\cap B)}{P(A)}. \]
C'est la probabilité que $B$ se réalise lorsque l'on sait que $A$ est réalisé. On la note aussi $P(B\,|\,A)$.
\end{definition}

\begin{propriete}[Règles de calcul]
$P_A$ est une probabilité : $0\leq P_A(B)\leq1$ et $P_A\big(\overline B\big)=1-P_A(B)$. De plus,
\[ P(A\cap B)=P(A)\times P_A(B)=P(B)\times P_B(A)\qquad\text{si }P(A)\neq0\text{ et }P(B)\neq0. \]
\end{propriete}

\essSection{2}{Arbres pondérés et probabilités totales}

\noindent
\begin{minipage}[c]{0.5\linewidth}
\centering
\begin{tikzpicture}[otArbre, level distance=27mm, level 1/.style={sibling distance=26mm},
  level 2/.style={sibling distance=13mm}]
  \node {}
    child {node {$A$}
      child {node (AB) {$B$} edge from parent node[otPoids,above,sloped] {$P_A(B)$}}
      child {node (AnB) {$\overline B$} edge from parent node[otPoids,below,sloped] {$P_A(\overline B)$}}
      edge from parent node[otPoids,above,sloped] {$P(A)$}}
    child {node {$\overline A$}
      child {node (nAB) {$B$} edge from parent node[otPoids,above,sloped] {$P_{\overline A}(B)$}}
      child {node (nAnB) {$\overline B$} edge from parent node[otPoids,below,sloped] {$P_{\overline A}(\overline B)$}}
      edge from parent node[otPoids,below,sloped] {$P(\overline A)$}};
  \node[otFeuille] at ([xshift=1.5mm]AB.east) {$P(A\cap B)$};
  \node[otFeuille] at ([xshift=1.5mm]AnB.east) {$P(A\cap\overline B)$};
  \node[otFeuille] at ([xshift=1.5mm]nAB.east) {$P(\overline A\cap B)$};
  \node[otFeuille] at ([xshift=1.5mm]nAnB.east) {$P(\overline A\cap\overline B)$};
\end{tikzpicture}\\[1mm]
\legende{Sur les branches du second niveau : des probabilités conditionnelles.}
\end{minipage}\hfill
\begin{minipage}[c]{0.46\linewidth}
\begin{theoreme}[Règles des arbres]
\raggedright
\begin{enumerate}[leftmargin=1.3em,itemsep=0.5mm,topsep=0.5mm,label={\color{otNavy}\arabic*.}]
  \item La somme des probabilités des branches issues d'un même nœud vaut $1$.
  \item La probabilité d'un chemin est le \textbf{produit} des probabilités de ses branches.
  \item La probabilité d'un événement est la \textbf{somme} des probabilités des chemins qui y mènent.
\end{enumerate}
\end{theoreme}
\end{minipage}

\begin{theoreme}[Formule des probabilités totales]
Si $A_1$, $A_2$, \dots, $A_n$ forment une \textbf{partition} de l'univers (événements deux à deux incompatibles, de réunion $\Omega$,
de probabilités non nulles), alors pour tout événement $B$ :
\[ P(B)=P(A_1\cap B)+\dots+P(A_n\cap B)=P(A_1)\,P_{A_1}(B)+\dots+P(A_n)\,P_{A_n}(B). \]
En particulier : $P(B)=P(A)\,P_A(B)+P\big(\overline A\big)\,P_{\overline A}(B)$.
\end{theoreme}

\begin{savoirfaire}[« Inverser » un conditionnement]
L'arbre donne $P_A(B)$ ; pour obtenir $P_B(A)$, on revient à la définition : $P_B(A)=\dfrac{P(A\cap B)}{P(B)}$, où $P(A\cap B)$
est lu sur un chemin et $P(B)$ calculé par la formule des probabilités totales. \emph{Exemple :} avec $P(A)=0{,}4$, $P_A(B)=0{,}3$
et $P_{\overline A}(B)=0{,}5$, on obtient $P(B)=0{,}12+0{,}3=0{,}42$ et $P_B(A)=\dfrac{0{,}12}{0{,}42}=\dfrac27$.
\end{savoirfaire}

\essSection{3}{Indépendance}

\begin{definition}[Événements indépendants]
Deux événements $A$ et $B$ sont \textbf{indépendants} si $P(A\cap B)=P(A)\times P(B)$. Lorsque $P(A)\neq0$, cela équivaut à
$P_A(B)=P(B)$ : savoir que $A$ est réalisé ne change pas la probabilité de $B$.
\end{definition}

\begin{propriete}[Indépendance et contraire]
Si $A$ et $B$ sont indépendants, alors $A$ et $\overline B$ le sont aussi (ainsi que $\overline A$ et $B$, $\overline A$ et $\overline B$).
\end{propriete}

\begin{demonstration}
$A$ est la réunion des événements incompatibles $A\cap B$ et $A\cap\overline B$, donc
$P\big(A\cap\overline B\big)=P(A)-P(A\cap B)=P(A)-P(A)P(B)=P(A)\big(1-P(B)\big)=P(A)\,P\big(\overline B\big)$.
\end{demonstration}

\begin{savoirfaire}[Succession d'épreuves indépendantes]
Lorsque des épreuves sont indépendantes (lancers successifs, tirages \textbf{avec remise}), la probabilité d'une issue
$(x_1\pv\dots\pv x_n)$ est le produit des probabilités de chaque résultat. Pour « au moins un succès » en $n$ épreuves de
probabilité de succès $p$, on passe au contraire : $1-(1-p)^{n}$.
\end{savoirfaire}

\begin{attention}[Pièges fréquents]
\begin{itemize}[label={\color{otOrange}\textbullet},leftmargin=1.4em,itemsep=0.6mm,topsep=0.5mm]
  \item $P_A(B)$ et $P_B(A)$ sont en général différents : un test « fiable à 99 \% » peut donner $P_T(M)$ très faible.
  \item Incompatibles ($A\cap B=\varnothing$) et indépendants sont deux notions différentes : deux événements incompatibles de
    probabilités non nulles ne sont \emph{jamais} indépendants.
  \item Sur un arbre, ne pas multiplier les probabilités de branches qui ne sont pas sur un même chemin.
  \item Tirages sans remise : les probabilités du second tirage dépendent du premier ; les épreuves ne sont pas indépendantes.
\end{itemize}
\end{attention}

%% =====================================================================
%%  EXERCICES
%% =====================================================================
\exercices
\rubrique{Découverte}

\begin{exercice}[1]{Arbres pondérés}
Dans chaque cas, construire un arbre pondéré, puis calculer la probabilité demandée.
\begin{questions}
  \item Une urne contient $3$ boules rouges et $2$ boules bleues. On tire successivement deux boules, sans remise. Probabilité
    d'obtenir deux boules rouges ?
  \item On lance deux fois une pièce équilibrée. Probabilité d'obtenir exactement un pile ?
  \item Une urne contient $4$ boules blanches et $6$ noires. On tire deux boules successivement, avec remise. Probabilité
    d'obtenir deux boules de couleurs différentes ?
  \item Un sac contient $5$ jetons : $2$ verts, $2$ jaunes et $1$ rouge. On tire un jeton, on le remet, puis on en tire un second.
    Probabilité d'obtenir deux jetons de la même couleur ?
  \item Une urne contient $4$ boules rouges et $6$ bleues. On tire successivement $3$ boules sans remise. Probabilité d'obtenir au
    moins une boule rouge ? \emph{On choisira la méthode la plus simple.}
\end{questions}
\end{exercice}

\begin{exercice}[1]{Calculs de base}
\begin{questions}
  \item On donne $P(A)=0{,}6$, $P(B)=0{,}4$ et $P(A\cap B)=0{,}2$. Calculer $P_B(A)$ et $P_A(B)$.
  \item On lance deux dés équilibrés. Sachant que la somme est supérieure à $9$, quelle est la probabilité que le premier dé
    montre $6$ ?
  \item Une urne contient $4$ boules rouges et $6$ bleues. On tire deux boules successivement, sans remise. Quelle est la
    probabilité que la deuxième soit rouge sachant que la première est rouge ?
  \item On donne $P_B(A)=0{,}7$ et $P(B)=0{,}3$. Calculer $P(A\cap B)$.
  \item On tire une carte d'un jeu de $52$ cartes. Sachant que c'est une figure (valet, dame ou roi), quelle est la probabilité
    que ce soit un cœur ?
\end{questions}
\end{exercice}

\rubrique{Entraînement}

\begin{exercice}[2]{Indépendance}
\begin{questions}
  \item On donne $P(A)=0{,}4$, $P(B)=0{,}5$ et $P(A\cap B)=0{,}2$. Les événements $A$ et $B$ sont-ils indépendants ?
  \item On lance deux dés équilibrés. Les événements « le premier dé montre $3$ » et « le second dé montre $5$ » sont-ils indépendants ?
  \item $A$ et $B$ sont indépendants, avec $P(A)=0{,}3$ et $P(B)=0{,}6$. Calculer $P(A\cup B)$.
  \item On tire une carte, on note sa couleur, on la remet, puis on tire une seconde carte. Les deux tirages sont-ils indépendants ?
  \item $A$ et $B$ sont indépendants, avec $P(A)=0{,}5$ et $P(B)=0{,}4$. Calculer $P\big(A\cap\overline B\big)$.
  \item On lance un dé équilibré ; on note $A$ : « obtenir un nombre pair » et $B$ : « obtenir $3$ ». Les événements $A$ et $B$
    sont-ils indépendants ? incompatibles ?
\end{questions}
\end{exercice}

\begin{exercice}[2]{Arbres pondérés et applications}
\begin{questions}
  \item Une urne contient $5$ boules rouges et $3$ vertes. On tire deux boules successivement sans remise. Probabilité d'obtenir
    une boule rouge puis une boule verte ?
  \item On lance deux fois un dé équilibré. Quelle est la probabilité d'avoir obtenu $6$ au premier lancer, sachant que la somme
    des deux résultats est $10$ ?
  \item Une classe de $30$ élèves compte $18$ filles, dont $10$ font du sport, et $12$ garçons, dont $8$ font du sport. On choisit
    un élève au hasard. À l'aide d'un arbre pondéré, calculer la probabilité qu'il fasse du sport.
  \item Un test de dépistage a une sensibilité de 98\,\% et une spécificité de 95\,\%, et 5\,\% de la population est atteinte.
    Quelle est la probabilité qu'un test soit positif ?
  \item On tire successivement deux cartes d'un jeu de $52$, sans remise. Probabilité que la seconde soit un as sachant que la
    première en est un ?
\end{questions}
\end{exercice}

\begin{exercice}[2]{Indépendance et applications}
\begin{questions}
  \item On lance deux dés équilibrés. Les événements « le premier dé montre un nombre pair » et « le second dé montre un nombre
    impair » sont-ils indépendants ?
  \item Une urne contient $6$ boules rouges et $4$ bleues. On tire successivement deux boules avec remise. Probabilité d'obtenir
    deux boules rouges ?
  \item On lance trois fois une pièce équilibrée. Probabilité d'obtenir exactement deux piles ?
  \item $A$ et $B$ sont indépendants, avec $P(A)=0{,}4$ et $P(B)=0{,}6$. Calculer $P(A\cup B)$.
  \item Un système électronique comporte deux composants qui fonctionnent indépendamment, avec les probabilités $0{,}9$ et $0{,}8$.
    Le système fonctionne si les deux composants fonctionnent. Quelle est la probabilité qu'il fonctionne ?
\end{questions}
\end{exercice}

\begin{exercice}[2]{Formule des probabilités totales}
\begin{questions}
  \item Une usine possède trois machines : $M_1$ fabrique 40\,\% de la production avec 2\,\% de pièces défectueuses, $M_2$ 35\,\%
    avec 3\,\% de défauts, et $M_3$ 25\,\% avec 5\,\% de défauts. Probabilité qu'une pièce prise au hasard soit défectueuse ?
  \item Dans une classe, 60\,\% des élèves sont des filles ; 70\,\% des filles et 50\,\% des garçons aiment les maths. Probabilité
    qu'un élève pris au hasard aime les maths ?
  \item Un test a une sensibilité de 95\,\% (il est positif pour 95\,\% des malades) et une spécificité de 90\,\% (il est négatif
    pour 90\,\% des personnes saines). La maladie touche 2\,\% de la population. Probabilité qu'un test soit positif ?
  \item Une boîte A contient $3$ boules rouges et $2$ bleues, une boîte B contient $1$ rouge et $4$ bleues. On choisit une boîte
    au hasard, puis on y tire une boule. Probabilité qu'elle soit rouge ?
  \item Un élève révise avec la probabilité $0{,}7$. S'il révise, il réussit avec la probabilité $0{,}9$ ; sinon, avec la
    probabilité $0{,}3$. Probabilité qu'il réussisse ?
\end{questions}
\end{exercice}

\begin{exercice}[2]{Inverser un conditionnement}
On reprend les situations de l'exercice 6.
\begin{questions}
  \item Une pièce est défectueuse. Quelle est la probabilité qu'elle provienne de la machine $M_3$ ?
  \item Le test d'une personne est positif. Quelle est la probabilité qu'elle soit réellement malade ? Commenter.
  \item On a tiré une boule rouge. Quelle est la probabilité qu'elle provienne de la boîte A ?
  \item L'élève a réussi. Quelle est la probabilité qu'il ait révisé ?
\end{questions}
\end{exercice}

\rubrique{Problèmes}

\begin{exercice}[3]{Tirages successifs sans remise}
Une urne contient $5$ boules rouges et $3$ boules bleues. On effectue $3$ tirages successifs sans remise.
\begin{questions}
  \item Quelle est la probabilité d'obtenir $3$ boules rouges ?
  \item Quelle est la probabilité d'obtenir exactement $2$ boules rouges ?
  \item Quelle est la probabilité d'obtenir au moins une boule bleue ?
  \item Sachant que la première boule tirée est rouge, quelle est la probabilité que les deux suivantes soient bleues ?
  \item Quelle est la probabilité que la troisième boule tirée soit rouge ? Comparer à la probabilité que la première le soit.
\end{questions}
\end{exercice}

\begin{exercice}[2]{Deux tirages avec remise}
Un sac contient $4$ jetons numérotés $1$, $2$, $3$, $4$. On tire un jeton, on note son numéro, on le remet, puis on en tire un second.
\begin{questions}
  \item Combien y a-t-il de résultats possibles ? Quelle est la probabilité d'obtenir deux numéros identiques ?
  \item Quelle est la probabilité que la somme des deux numéros soit égale à $5$ ?
  \item Quelle est la probabilité que le produit des deux numéros soit pair ?
  \item Sachant que la somme est paire, quelle est la probabilité que les deux numéros soient pairs ?
\end{questions}
\end{exercice}

\begin{exercice}[3]{Applications complexes}
\begin{questions}
  \item Une usine produit des pièces sur trois chaînes : A (50\,\% des pièces, 1\,\% de défauts), B (30\,\%, 2\,\% de défauts) et
    C (20\,\%, 4\,\% de défauts). Calculer la probabilité qu'une pièce soit défectueuse, puis la probabilité qu'une pièce
    défectueuse provienne de la chaîne C.
  \item On tire successivement trois cartes d'un jeu de $52$, sans remise. Probabilité d'obtenir exactement deux as ?
  \item Un test a une sensibilité de 96\,\% et une spécificité de 94\,\% ; la maladie touche 3\,\% de la population. Calculer la
    probabilité qu'un test soit positif, puis celle d'être malade sachant que le test est positif.
  \item Une urne contient $5$ boules rouges et $3$ bleues. On tire successivement $3$ boules \emph{avec} remise. Probabilité
    d'obtenir au moins une boule bleue ? Comparer avec le résultat sans remise (exercice 8).
  \item On lance quatre fois un dé équilibré. Probabilité d'obtenir au moins un $6$ ?
\end{questions}
\end{exercice}

\begin{exercice}[3]{Devoir de synthèse}
\begin{questions}
  \item On donne $P(A)=0{,}5$, $P(B)=0{,}4$ et $P(A\cap B)=0{,}15$. Calculer $P_B(A)$ et $P_A(B)$.
  \item On tire une carte d'un jeu de $52$. Sachant que c'est un cœur, quelle est la probabilité que ce soit un as ?
  \item Une urne contient $4$ boules rouges et $6$ bleues. On tire deux boules successivement sans remise. Probabilité d'obtenir
    deux boules de couleurs différentes ?
  \item On lance trois fois une pièce équilibrée. Probabilité d'obtenir exactement deux faces ?
  \item On lance deux dés. Les événements « le premier dé montre un nombre pair » et « le second dé montre un nombre supérieur
    à $3$ » sont-ils indépendants ?
  \item $A$ et $B$ sont indépendants, avec $P(A)=0{,}6$ et $P(B)=0{,}5$. Calculer $P(A\cup B)$.
  \item Une usine a deux machines : $M_1$ (60\,\% de la production, 1\,\% de défauts) et $M_2$ (40\,\%, 3\,\% de défauts). Calculer
    la probabilité qu'un produit soit défectueux, puis celle qu'un produit défectueux provienne de $M_2$.
  \item Dans une classe, 55\,\% des élèves sont des filles ; 75\,\% des filles et 60\,\% des garçons aiment les sciences.
    Probabilité qu'un élève aime les sciences ?
  \item Un test a une sensibilité de 97\,\% et une spécificité de 92\,\% ; 4\,\% de la population est malade. Probabilité d'être
    malade sachant que le test est positif ?
\end{questions}
\end{exercice}

\begin{exercice}[3]{Le paradoxe des tests de dépistage}
Une maladie touche une proportion $p$ de la population. Un test de dépistage a une sensibilité de 99\,\%, soit $P_M(T)=0{,}99$,
et une spécificité de 98\,\%, soit $P_{\overline M}\big(\overline T\big)=0{,}98$, où $M$ est l'événement « la personne est malade » et
$T$ l'événement « le test est positif ».
\partieExo{A}{Une maladie rare : $p=0{,}01$}
\begin{questions}
  \item Compléter un tableau d'effectifs pour une population de $10\,000$ personnes (malades ou non, test positif ou négatif).
  \item En déduire $P(T)$, puis la probabilité qu'une personne dont le test est positif soit malade. Commenter.
  \item On fait passer un second test à toute personne testée positive. On admet que, chez les malades comme chez les personnes
    saines, les résultats des deux tests sont indépendants. Quelle est la probabilité qu'une personne dont les deux tests sont
    positifs soit malade ?
\end{questions}
\partieExo{B}{Influence de la proportion de malades}
\begin{questions}[resume]
  \item Montrer que $P(T)=0{,}97\,p+0{,}02$, puis que $P_T(M)=f(p)$, où $f$ est la fonction définie sur $\intOF{0}{1}$ par
    $f(p)=\dfrac{99\,p}{97\,p+2}$.
  \item Étudier les variations de $f$ sur $\intOF{0}{1}$.
  \item À partir de quelle proportion de malades la probabilité $P_T(M)$ est-elle au moins égale à $0{,}9$ ?
\end{questions}
\end{exercice}

\begin{exercice}[3]{Une suite de probabilités}
Chaque matin, Léa se rend au lycée à vélo ou en bus. Si elle a pris le vélo un jour, elle le reprend le lendemain avec la
probabilité $0{,}7$ ; si elle a pris le bus, elle prend le vélo le lendemain avec la probabilité $0{,}4$. Le premier jour, elle
prend le vélo. Pour $n\geq1$, on note $V_n$ l'événement « Léa prend le vélo le jour $n$ » et $p_n=P(V_n)$ ; ainsi $p_1=1$.
\begin{questions}
  \item Calculer $p_2$ et $p_3$.
  \item Représenter par un arbre pondéré la situation entre le jour $n$ et le jour $n+1$, puis montrer que $p_{n+1}=0{,}3\,p_n+0{,}4$.
  \item On pose $u_n=p_n-\dfrac47$. Montrer que la suite $(u_n)$ est géométrique, puis exprimer $p_n$ en fonction de $n$.
  \item Déterminer la limite de $(p_n)$ et l'interpréter.
  \item Que renvoie la fonction Python ci-dessous ? Interpréter.
\begin{python}
def jour():
    p = 1
    n = 1
    while p - 4/7 >= 0.001:
        p = 0.3 * p + 0.4
        n = n + 1
    return n
\end{python}
\end{questions}
\end{exercice}

\begin{exercice}[3]{Fiabilité d'un système}
Un appareil comporte des composants identiques, qui fonctionnent indépendamment les uns des autres : chacun fonctionne pendant
un an avec la probabilité $0{,}9$. Des composants montés en \textbf{série} assurent le fonctionnement si \emph{tous} fonctionnent ;
montés en \textbf{parallèle}, si \emph{au moins un} fonctionne.
\begin{questions}
  \item Quelle est la probabilité que deux composants montés en série fonctionnent ? Avec $n$ composants en série, à partir de
    quelle valeur de $n$ cette probabilité devient-elle inférieure à $0{,}5$ ?
  \item Même question pour deux composants montés en parallèle. Combien faut-il monter de composants en parallèle pour que la
    probabilité de fonctionnement soit au moins égale à $0{,}999\,9$ ?
  \item On dispose de quatre composants. Le montage (a) est formé de deux branches en parallèle, chacune comportant deux composants
    en série ; le montage (b) est formé de deux blocs en série, chacun comportant deux composants en parallèle. Quel est le
    montage le plus fiable ?
  \item Pour deux composants montés en parallèle, sachant que l'appareil fonctionne, quelle est la probabilité que les deux
    composants fonctionnent ?
\end{questions}
\begin{center}
\begin{tikzpicture}[x=1mm,y=1mm,
  comp/.style={draw=otNavy,fill=otNavyTint,minimum width=7mm,minimum height=4.5mm,inner sep=0pt,font=\scriptsize}]
  \node[font=\scriptsize,text=otMuted] at (-6,9) {(a)};
  \draw[otMuted,line width=0.6pt] (-4,0) -- (0,0) (0,5) -- (0,-5) (0,5) -- (4,5) (0,-5) -- (4,-5)
    (11,5) -- (15,5) (11,-5) -- (15,-5) (22,5) -- (26,5) (22,-5) -- (26,-5) (26,5) -- (26,-5) (26,0) -- (30,0);
  \node[comp] at (7.5,5) {C}; \node[comp] at (18.5,5) {C}; \node[comp] at (7.5,-5) {C}; \node[comp] at (18.5,-5) {C};
  \begin{scope}[xshift=48mm]
  \node[font=\scriptsize,text=otMuted] at (-6,9) {(b)};
  \draw[otMuted,line width=0.6pt] (-4,0) -- (0,0) (0,5) -- (0,-5) (0,5) -- (4,5) (0,-5) -- (4,-5) (11,5) -- (15,5) (11,-5) -- (15,-5)
    (15,5) -- (15,-5) (15,0) -- (19,0) (19,5) -- (19,-5) (19,5) -- (23,5) (19,-5) -- (23,-5) (30,5) -- (34,5) (30,-5) -- (34,-5)
    (34,5) -- (34,-5) (34,0) -- (38,0);
  \node[comp] at (7.5,5) {C}; \node[comp] at (7.5,-5) {C}; \node[comp] at (26.5,5) {C}; \node[comp] at (26.5,-5) {C};
  \end{scope}
\end{tikzpicture}
\end{center}
\end{exercice}
